\providecommand{\divby}{\par\smallskip\noindent\hrulefill\par\smallskip} % разделитель частей; если в стиле он есть, возьмётся стилевой

\ListokName{18. Вероятность}

% Шапка: набрана вручную, потому что команды шапки для уровней α и ℵ нам неизвестны.
% Если в стиле есть своя (как \shapkaUstnoA для уровня А), замените ею эти три строки.
\noindent\rlap{\makebox[\linewidth]{\itshape 2026/2027. Школа 179. КЛ 25-29}}\hfill{\LARGE$\alpha$}\par\medskip
\centerline{\rule{0.62\linewidth}{0.4pt}}\par\smallskip
\centerline{\large\bfseries\scshape 18. Вероятность}\par\medskip

Бросим симметричную монету тысячу раз. Предсказать, как она упадёт в очередной раз, нельзя, но орлов, скорее всего, будет примерно половина. Бросок монеты~— \textit{эксперимент}: процедура, которую можно повторять сколько угодно раз, а результат от раза к разу разный. \textit{Событие}~— утверждение о результате эксперимента: после каждого опыта можно сказать, произошло оно или нет: «выпал орёл», «на кубике выпало чётное число». Если в $N$ повторениях эксперимента событие произошло $k$ раз, число $k/N$ называют \textit{частотой} события.

У симметричного кубика грани ничем не отличаются друг от друга, поэтому в длинной серии бросков каждая грань обычно выпадает примерно в шестой части случаев. Слова «случайный выбор» («тянут наугад», «садится на случайное место») означают то же самое: в длинной серии опытов каждый вариант выбора встречается примерно одинаково часто.

\zp \unskip$^{\circ}$\ На гранях симметричной двадцатигранной кости написаны числа от 1 до 20. Какова примерная частота выпадения \textbf{а)} простого числа; \textbf{б)} числа, делящегося на 3 или на 5?

\zp Симметричный кубик бросают; если выпала шестёрка, бросок не засчитывают и кубик перебрасывают. Какова примерная частота единицы среди засчитанных бросков? А чётного числа?

Если эксперимент повторять много раз, частота события колеблется около одного и того же числа~— где бы и кем бы эксперимент ни проводился. Это число называют \textit{вероятностью} события $A$ и обозначают $P(A)$. Это не теорема, а наблюдение о мире; на нём основана вера в то, что теорию вероятностей можно применять к реальной жизни. Слово «вероятность» используют и в других смыслах: например, говорят о вероятности единичного события, которое нельзя повторить много раз,~— что определённая команда выиграет завтрашний матч. В этом листке мы говорим только о вероятности как о частоте.

Событие «произошло хотя бы одно из событий $A$ и $B$» обозначают $A\cup B$, «произошли и $A$, и $B$»~— $A\cap B$, «$A$ не произошло»~— $\overline{A}$; его называют \textit{противоположным} к $A$. События $A$ и $B$ \textit{несовместны}, если они не могут произойти одновременно.

\zp \unskip$^{\circ}$\ \textbf{а)} Объясните, почему $P(A\cup B)=P(A)+P(B)$, если события $A$ и $B$ несовместны. \textbf{б)} Чему равна вероятность события, которое происходит при любом результате эксперимента? \textbf{в)} Объясните, почему $P(A)+P(\overline{A})=1$.

\zp \unskip$^{\circ}$\ Верно ли, что $P(A\cup B)=P(A)+P(B)$ для любых событий $A$ и $B$? Если нет, сформулируйте верное равенство.

\divby

\zp \unskip$^{\circ}$\ Две симметричные монеты бросают одновременно. Найдите вероятность того, что выпадут два орла; ровно один орёл.

\z Бросают два симметричных кубика, красный и синий, и складывают выпавшие числа.
\lett \unskip$^{\circ}$\ Какая сумма имеет большую вероятность: 3 или 10? Во сколько раз?
\lett Какая сумма имеет наибольшую вероятность?

\zp \unskip$^{\circ}$\ Десять школьников по очереди тянут билеты из десяти; вытянутый билет на стол не возвращают. Вася выучил один билет. Что ему выгоднее: тянуть первым или последним?

\z Восемь команд разной силы играют турнир по олимпийской системе: три круга, проигравший выбывает, в каждом матче побеждает более сильная команда. Пары первого круга определяет жеребьёвка. С какой вероятностью \leth в финале встретятся две сильнейшие команды; \leth в финал выйдет третья по силе команда?

\z За круглый стол в случайном порядке садятся 10 мальчиков и 10 девочек, среди них Маша. С какой вероятностью \leth \unskip$^{\circ}$\ справа от Маши сидит мальчик; \leth оба её соседа~— мальчики?

\z Жук сидит в вершине проволочного каркаса куба. Каждую минуту он переползает по ребру в одну из трёх соседних вершин, выбирая её случайно. С какой вероятностью
\lett через 3 минуты он окажется в вершине, противоположной исходной;
\lett через 4 минуты он снова окажется в исходной вершине?

\z С какой вероятностью ни разу не выпадут два орла подряд, если монету бросают \leth 4 раза; \leth 10 раз?

\z На выборах бюллетени вынимают из урны по одному в случайном порядке и сразу подсчитывают. С какой вероятностью в какой-то момент подсчёта у кандидатов А и Б окажется поровну голосов, если за А подано \leth 3 голоса, а за Б~— 2; \leth \unskip$^{\star}$\ 30 голосов, а за Б~— 20?

\z Каждое из чисел 1, 2, \ldots, 8 красят в красный или синий цвет, бросая для каждого числа симметричную монету.
\lett \unskip$^{\circ}$\ Что вероятнее: сумма красных чисел больше суммы синих или меньше?
\lett Найдите вероятность того, что сумма красных чисел больше.

\zp \unskip$^{\circ}$\ В урне 3 белых и 2 чёрных шара. Их вынимают наугад по одному, пока не появится белый. С какой вероятностью придётся вынуть не меньше трёх шаров?

\z Игра состоит из партий; в каждой партии каждая из сторон выигрывает с вероятностью $1/2$, ничьих не бывает. Если игру прервали, приз делят в отношении вероятностей победы, которые были у сторон в момент остановки.
\lett \unskip$^{\circ}$\ Двое играли до трёх выигранных партий и прервали игру при счёте 2:1. Как разделить приз?\footnote{Задачу о разделе приза в прерванной игре обсуждали в 1654 году в переписке Блез Паскаль и Пьер Ферма. Эту переписку считают началом теории вероятностей.}
\lett Команды Знатоков и Телезрителей играют до шести побед. Сейчас счёт 2:4 в пользу Телезрителей. С какой вероятностью победят Телезрители?

\zp Аня бросает $n+1$ симметричную монету, а Боря~— $n$ монет. С какой вероятностью у Ани выпадет больше орлов, чем у Бори?

\divby

\textit{Вероятностная модель} эксперимента~— это конечное множество $\Omega$ и числа $p(\omega)\ge 0$, по одному для каждого $\omega\in\Omega$, в сумме равные 1. Элементы $\Omega$~— \textit{элементарные исходы}: несовместные события, из которых при каждом проведении эксперимента происходит ровно одно; $p(\omega)$~— их вероятности. Событию отвечает подмножество $A\subset\Omega$, и $P(A)$~— сумма $p(\omega)$ по всем $\omega\in A$. Один эксперимент можно описать разными моделями; если обе правильно описывают эксперимент, вероятность события в них одна и та же.

\zp \unskip$^{\dag}$\ Докажите, исходя из определения модели: \textbf{а)} $P(\Omega)=1$ и $P(\emptyset)=0$; \textbf{б)} если $A\subset B$, то $P(A)\le P(B)$; \textbf{в)} $P(\overline{A})=1-P(A)$; \textbf{г)} $P(A\cup B)=P(A)+P(B)-P(A\cap B)$.

\zp \textbf{а)} Опишите вероятностные модели экспериментов из задач 6 и 10: что такое $\Omega$ и чему равны $p(\omega)$? \textbf{б)} Для задачи 9 постройте две модели: в первой элементарный исход~— рассадка всех двадцати человек, во второй~— пара «кто сидит слева от Маши, кто справа». Сколько исходов в каждой модели? Найдите в обеих ответ на вопрос 9б.

\zp \unskip$^{\circ}$\ Бросают две одинаковые монеты. Петя рассуждает так: исходов три~— два орла, две решки, орёл и решка,~— значит, монеты упадут по-разному с вероятностью $1/3$. Вася насчитал четыре исхода и получил $1/2$. Кто прав? В чём ошибка другого?

\zp Десять человек, среди них Петя и Вася, садятся в случайном порядке \textbf{а)} в ряд; \textbf{б)} за круглый стол. С какой вероятностью между Петей и Васей сидят ровно $k$ человек (за столом~— с одной из сторон)?

\zp За круглый стол в случайном порядке садятся 4 мальчика и 4 девочки. С какой вероятностью их можно разбить на 4 пары соседей так, чтобы в каждой паре были мальчик и девочка?

\zp \unskip$^{\star}$\ Каждый из $n$ пассажиров купил билет на $n$-местный самолёт. Первой вошла рассеянная старушка и села на случайное место. Каждый следующий пассажир садится на своё место, если оно свободно, а если занято~— на случайное свободное. С какой вероятностью последний пассажир сядет на своё место?
